% =====================================================================
% Identidade visual de gráficos e tabelas
%
% Paleta e macros adaptadas do código-fonte (LaTeX) de dois artigos:
%   - OA-RAG: Ontology-based Agentic RAG (arXiv:2610.36335) -- paleta,
%     destaques de tabela (\tfirst, \tsecond, \gain, \loss, \cmark,
%     \nomark, linha/coluna em destaque), legenda com "chips" e o estilo
%     de barras horizontais com rótulo interno;
%   - Hossain e Sarıyüce, Core-based Hierarchies for Efficient GraphRAG
%     (arXiv:2603.05207) -- macro \heatmap de células de tabela.
% =====================================================================

\usepackage{colortbl}
\usepackage{pifont}
\usepackage{fontawesome5}   % ícones da figura de arquitetura
\usepackage{placeins}        % \FloatBarrier: mantém figuras dentro da própria seção
\usetikzlibrary{patterns, shapes.arrows}

% --- Paleta principal (OA-RAG) -----------------------------------------
\definecolor{vAzul}{HTML}{2A78D6}
\definecolor{vVerde}{HTML}{1BAF7A}
\definecolor{vLaranja}{HTML}{EB6834}
\definecolor{vAmbar}{HTML}{EDA100}
\definecolor{vRoxo}{HTML}{4A3AA7}
\definecolor{vVermelho}{HTML}{E34948}
\definecolor{vTeal}{HTML}{2A9D8F}
\definecolor{vCinza}{HTML}{9B9A94}
\definecolor{vNeutro}{HTML}{DDDCD5}

% Cada modo de busca tem uma cor fixa em todo o documento
\colorlet{modoBasic}{vAzul}
\colorlet{modoLocal}{vVerde}
\colorlet{modoGlobal}{vRoxo}
\colorlet{modoDrift}{vLaranja}

% --- Tabelas (OA-RAG) --------------------------------------------------
\colorlet{linhaDestaque}{black!6}
\colorlet{colunaDestaque}{black!8}
\newcommand{\tfirst}[1]{{\color{red!80!black}#1}}
\newcommand{\tsecond}[1]{{\color{blue!80!black}#1}}
\newcommand{\gain}[1]{\;{\color{green!50!black}\ensuremath{\uparrow}\,#1}}
\newcommand{\loss}[1]{\;{\color{red!75!black}\ensuremath{\downarrow}\,#1}}
\newcommand{\cmark}{\ding{51}}
\newcommand{\nomark}{{\color{black!35}--}}
\newcommand{\tcok}{\textcolor{green!50!black}{\ding{51}}}
\newcommand{\tcno}{\textcolor{red!75!black}{\ding{55}}}

% Amostra de cor para legendas (o "\chip" do OA-RAG)
\newcommand{\chip}[1]{\tikz[baseline=-0.6ex]\filldraw[fill=#1, draw=black!40, line width=0.3pt] (0,-0.09) rectangle (0.30,0.09);\hspace{2.5pt}}

% Barra proporcional dentro de célula de tabela (valor em [0,1]), no
% estilo das células de acurácia do OA-RAG: \barcell{cor}{valor}{rótulo}
\newcommand{\barcell}[3]{%
  \tikz[baseline=-0.6ex]{%
    \fill[black!5] (0,-0.11) rectangle (1.6,0.11);%
    \fill[#1] (0,-0.11) rectangle ({1.6*#2},0.11);%
    \node[anchor=west, inner sep=0pt, font=\footnotesize] at ({1.6*#2+0.06},0) {#3};%
  }%
}

% --- Posicionamento de figuras (OA-RAG) ---------------------------------
% Limites padrão do LaTeX deixam figuras altas presas na fila de floats,
% empurrando todas as seguintes para o fim do capítulo. Mesmos valores do
% preâmbulo do OA-RAG.
\renewcommand{\topfraction}{.92}\renewcommand{\dbltopfraction}{.92}
\renewcommand{\textfraction}{.06}\renewcommand{\floatpagefraction}{.72}

% --- Heatmap de células (Hossain e Sarıyüce) ----------------------------
% Branco (0) -> amarelo (50) -> vermelho (100); imprime o valor.
\newcommand{\heatmap}[1]{%
    \pgfmathsetmacro{\val}{#1}%
    \pgfmathsetmacro{\r}{ifthenelse(\val<50, \val/50, 1)}%
    \pgfmathsetmacro{\g}{ifthenelse(\val<50, 1, 2-\val/50)}%
    \pgfmathsetmacro{\b}{0}%
    \edef\temp{\noexpand\cellcolor[rgb]{\r,\g,\b}}%
    \temp #1%
}
